\documentclass[UTF8,zihao=-4]{ctexart}
\usepackage[a4paper,margin=2.25cm]{geometry}
\usepackage{amsmath,amssymb,bm}
\usepackage{booktabs,longtable,array,multirow}
\usepackage{tikz}
\usetikzlibrary{arrows.meta,positioning,shapes.geometric,fit,calc}
\usepackage{graphicx}
\usepackage{xcolor}
\usepackage{enumitem}
\usepackage{hyperref}
\usepackage{tcolorbox}
\usepackage{listings}
\usepackage{caption}
\usepackage{float}
\usepackage{titlesec}
\hypersetup{colorlinks=true,linkcolor=blue!55!black,urlcolor=blue!55!black,citecolor=blue!55!black}
\definecolor{mainblue}{RGB}{29,63,112}
\definecolor{lightblue}{RGB}{237,244,255}
\definecolor{lightgray}{RGB}{248,248,248}
\definecolor{codebg}{RGB}{248,248,248}
\titleformat{\section}{\Large\bfseries\color{mainblue}}{\thesection}{0.8em}{}
\titleformat{\subsection}{\large\bfseries\color{mainblue}}{\thesubsection}{0.8em}{}
\setlist[itemize]{leftmargin=2em,itemsep=0.25em,topsep=0.3em}
\setlist[enumerate]{leftmargin=2em,itemsep=0.25em,topsep=0.3em}
\tcbset{colback=lightblue,colframe=mainblue,arc=1mm,boxrule=0.6pt,left=1.2mm,right=1.2mm,top=1mm,bottom=1mm}
\lstset{basicstyle=\ttfamily\small,backgroundcolor=\color{codebg},frame=single,breaklines=true,columns=fullflexible,keepspaces=true,keywordstyle=\color{blue!60!black}\bfseries,commentstyle=\color{green!35!black},stringstyle=\color{red!50!black},showstringspaces=false}

\title{TOPSIS、模糊综合评价与熵权法的共同建模思想\\\large 从数据矩阵到综合评价结论的统一套路}
\author{数学建模方法总结}
\date{\today}

\begin{document}
\maketitle

\begin{abstract}
TOPSIS、模糊综合评价法（FCE）和熵权法看起来公式不同，但都属于多指标综合评价或多准则决策问题中的常用工具。它们的共同主线可以概括为：先把现实问题整理成“评价对象—评价指标”的数据矩阵，再通过正向化、标准化或归一化消除指标方向和量纲差异，接着确定指标权重，最后把多维指标信息合成为可以比较的排序、等级或综合得分。本文按建模论文写法总结三类方法的相似点、主旨、数据处理逻辑、权重逻辑、组合模型、结论评价模板、可拓展方向和 Python 代码。重点回答：为什么要找最大值、为什么要建矩阵、标准化与归一化有什么区别、TOPSIS/FCE/熵权法分别在同一条评价链条中承担什么角色。
\end{abstract}

\tableofcontents
\newpage

\section{一句话总概括}
\begin{tcolorbox}
这些方法的主旨不是“套公式”，而是把一个多指标、多方案、多尺度的问题，转化为一个可计算的评价矩阵；再通过同向化、无量纲化、赋权和综合合成，把原来不能直接比较的数据压缩成一个有解释性的排序、等级或得分。
\end{tcolorbox}

更简单地说，它们共同解决的是：
\begin{itemize}
  \item 指标很多，不能只看一个数；
  \item 指标单位不同，例如万元、百分比、人数、浓度、时间；
  \item 指标方向不同，有的越大越好，有的越小越好，有的越接近某目标越好；
  \item 指标重要性不同，需要权重；
  \item 最后必须给出“哪个更好”“属于哪个等级”“改进重点在哪里”。
\end{itemize}

可以把共同套路写成八个字：\textbf{矩阵化、同向化、无量纲化、赋权、合成、评价}。若写成完整流程，则是：

\begin{center}
\begin{tikzpicture}[node distance=9mm, every node/.style={font=\small}, box/.style={rectangle, rounded corners, draw=mainblue, fill=lightblue, minimum width=3.2cm, minimum height=8mm, align=center}, arrow/.style={-Latex, thick, draw=mainblue}]
\node[box] (a) {现实评价问题\\对象与指标};
\node[box, right=of a] (b) {原始数据矩阵\\$X=(x_{ij})$};
\node[box, right=of b] (c) {正向化\\统一“越大越好”};
\node[box, below=of c] (d) {标准化/归一化\\消除量纲};
\node[box, left=of d] (e) {确定权重\\AHP/熵权/组合};
\node[box, left=of e] (f) {综合合成\\TOPSIS/FCE/加权得分};
\node[box, below=of f] (g) {排序、分级\\结论与建议};
\draw[arrow] (a)--(b); \draw[arrow] (b)--(c); \draw[arrow] (c)--(d); \draw[arrow] (d)--(e); \draw[arrow] (e)--(f); \draw[arrow] (f)--(g);
\end{tikzpicture}
\end{center}

\section{共同的数据结构：为什么都要建立矩阵}
设有 $n$ 个评价对象和 $m$ 个评价指标：
\[
A=\{A_1,A_2,\ldots,A_n\},\qquad C=\{C_1,C_2,\ldots,C_m\}.
\]
原始数据矩阵为
\[
X=(x_{ij})_{n\times m}=\begin{bmatrix}
x_{11}&x_{12}&\cdots&x_{1m}\\
x_{21}&x_{22}&\cdots&x_{2m}\\
\vdots&\vdots&\ddots&\vdots\\
x_{n1}&x_{n2}&\cdots&x_{nm}
\end{bmatrix},
\]
其中 $x_{ij}$ 表示第 $i$ 个对象在第 $j$ 个指标上的取值。

\subsection{矩阵的意义}
矩阵不是形式主义，而是把复杂评价问题整理成统一计算对象：
\begin{itemize}
  \item \textbf{行}表示评价对象，例如城市、企业、方案、月份、样本；
  \item \textbf{列}表示评价指标，例如成本、收益、风险、满意度、污染物浓度；
  \item 每个元素 $x_{ij}$ 是一个可计算的观测值、评分或统计量。
\end{itemize}

TOPSIS 用矩阵寻找正理想解和负理想解；熵权法用矩阵计算各指标的信息差异；FCE 用矩阵表达各指标对各评语等级的隶属度。因此，矩阵是这些方法的共同起点。

\section{共同的数据处理：正向化、标准化、归一化与距离打分}
\subsection{第一步：正向化处理}
正向化的目标是把所有指标统一为“数值越大，评价越好”。常见指标类型如下。

\begin{longtable}{p{2.8cm}p{5.2cm}p{5.4cm}}
\toprule
指标类型 & 含义 & 常见处理方式 \\
\midrule
最大型 & 越大越好，如收入、合格率、绿化率 & 保持不变：$x'_{ij}=x_{ij}$ \\
最小型 & 越小越好，如成本、污染、故障率 & $x'_{ij}=\max_i x_{ij}-x_{ij}$ 或倒数变换 \\
中间型 & 越接近目标值越好，如适宜温度 & $x'_{ij}=1-\frac{|x_{ij}-a_j|}{\max_i |x_{ij}-a_j|}$ \\
区间型 & 落在合理区间最好，如合理负债率 & 先计算偏离区间的距离，再转化为越大越好 \\
\bottomrule
\end{longtable}

\subsection{第二步：标准化处理}
标准化的核心作用是\textbf{消除量纲和数量级影响}。例如“GDP 为 10000 亿元”和“满意度为 0.8”不能直接相加。常用极差标准化为
\[
r_{ij}=\frac{x'_{ij}-\min_i x'_{ij}}{\max_i x'_{ij}-\min_i x'_{ij}}.
\]
该公式能使 $r_{ij}\in[0,1]$，但一般\textbf{不能保证同一列或同一行的和为 1}。

\subsection{第三步：归一化处理}
归一化在不同上下文中含义略有差别。熵权法常用比例归一化：
\[
p_{ij}=\frac{r_{ij}}{\sum_{i=1}^n r_{ij}},\qquad \sum_{i=1}^n p_{ij}=1.
\]
TOPSIS 常用向量归一化：
\[
z_{ij}=\frac{x'_{ij}}{\sqrt{\sum_{i=1}^n (x'_{ij})^2}}.
\]
所以要注意：\textbf{标准化通常是为了去量纲；归一化常常还要求比例和为 1 或向量长度为 1。}

\subsection{距离打分与样本最值}
在 TOPSIS 或一些距离评分中，最大值和最小值通常取当前样本中的最大值和最小值，而不一定是理论最大值和理论最小值。原因是很多指标没有稳定的理论边界，例如 GDP 增速、创新能力指数、综合满意度等。此时取样本最值，本质上是在当前评价对象集合内部建立相对比较标准。

\begin{tcolorbox}
标准化、归一化和距离打分的关系可以这样记：标准化解决“能不能相加比较”；归一化解决“能不能看成比例或单位向量”；距离打分解决“离最好有多近、离最差有多远”。
\end{tcolorbox}

\section{三类方法的主旨与位置}
\subsection{TOPSIS：用“接近最好、远离最差”排序}
TOPSIS 的基本思想是：最优方案应当尽可能接近正理想解，同时尽可能远离负理想解。设标准化或加权标准化矩阵为 $Z=(z_{ij})$，则正理想解和负理想解为
\[
Z^+=(Z^+_1,\ldots,Z^+_m),\qquad Z^-=(Z^-_1,\ldots,Z^-_m),
\]
其中
\[
Z^+_j=\max_i z_{ij},\qquad Z^-_j=\min_i z_{ij}.
\]
第 $i$ 个对象到正理想解和负理想解的距离为
\[
D_i^+=\sqrt{\sum_{j=1}^m (z_{ij}-Z_j^+)^2},\qquad
D_i^-=\sqrt{\sum_{j=1}^m (z_{ij}-Z_j^-)^2}.
\]
相对贴近度为
\[
S_i=\frac{D_i^-}{D_i^+ + D_i^-},\qquad 0\le S_i\le 1.
\]
$S_i$ 越大，说明该对象越接近理想方案、越远离负理想方案，综合排序越靠前。

\subsection{FCE：用“属于各等级的程度”分级}
模糊综合评价法的核心不是直接给一个硬分数，而是给出评价对象对若干评语等级的隶属程度。设因素集为
\[
U=\{u_1,u_2,\ldots,u_m\},
\]
评语集为
\[
V=\{v_1,v_2,\ldots,v_s\},
\]
权重向量为
\[
W=(w_1,w_2,\ldots,w_m),\qquad \sum_{j=1}^m w_j=1.
\]
隶属度矩阵为
\[
R=\begin{bmatrix}
r_{11}&r_{12}&\cdots&r_{1s}\\
r_{21}&r_{22}&\cdots&r_{2s}\\
\vdots&\vdots&\ddots&\vdots\\
r_{m1}&r_{m2}&\cdots&r_{ms}
\end{bmatrix}.
\]
综合评价向量为
\[
B=WR=(b_1,b_2,\ldots,b_s).
\]
若 $b_k$ 最大，则可按最大隶属度原则判断综合等级为 $v_k$。若评语集有顺序，还可以给评语赋分 $Q=(q_1,\ldots,q_s)$，计算综合得分
\[
S=BQ^T.
\]

\subsection{熵权法：用“数据差异大小”客观赋权}
熵权法本身主要是\textbf{赋权方法}，不是完整的排序或分级模型。它的思想是：某指标在不同评价对象之间差异越明显，越能区分对象，权重越高；若某指标几乎所有对象都差不多，则区分作用弱，权重低。

将标准化矩阵转为比例矩阵
\[
p_{ij}=\frac{r_{ij}}{\sum_{i=1}^n r_{ij}}.
\]
第 $j$ 个指标的信息熵为
\[
e_j=-\frac{1}{\ln n}\sum_{i=1}^n p_{ij}\ln p_{ij}.
\]
差异系数为
\[
g_j=1-e_j.
\]
熵权为
\[
w_j=\frac{g_j}{\sum_{j=1}^m g_j}.
\]
注意：熵权高不等于现实中“绝对重要”，而是表示该指标在当前样本中更能拉开差距。

\section{三者的相似点与不同点}
\begin{longtable}{p{3cm}p{4.1cm}p{4.1cm}p{4.1cm}}
\toprule
比较角度 & TOPSIS & FCE & 熵权法 \\
\midrule
共同问题 & \multicolumn{3}{p{12.3cm}}{都用于多指标综合评价或多准则决策，要求先建立对象—指标矩阵，并处理指标方向、量纲和权重。} \\
核心问题 & 谁离理想方案更近 & 对象属于哪个等级的程度更高 & 哪些指标在样本中区分力更强 \\
主要输出 & 贴近度、排序 & 隶属度向量、等级、得分 & 指标权重 \\
是否直接排序 & 是 & 可通过赋分排序 & 通常需与加权求和、TOPSIS 或 FCE 结合 \\
适合场景 & 多方案优劣排序 & 等级边界模糊、专家评价、定性定量混合 & 客观数据较完整、希望减少主观赋权 \\
关键矩阵 & 加权标准化决策矩阵 & 隶属度矩阵 & 标准化矩阵、比例矩阵 \\
核心公式 & $S_i=D_i^-/(D_i^+ + D_i^-)$ & $B=WR$ & $w_j=g_j/\sum g_j$ \\
常见组合 & 熵权TOPSIS、模糊TOPSIS & AHP-FCE、熵权FCE & 熵权TOPSIS、熵权FCE、AHP-熵权组合 \\
\bottomrule
\end{longtable}

\section{它们可以怎么组合}
\subsection{AHP-TOPSIS}
AHP 用专家判断矩阵确定权重，TOPSIS 用距离贴近度排序。适合专家经验明确、指标数据也比较完整的方案选择问题。

\subsection{熵权-TOPSIS}
熵权法客观确定权重，TOPSIS 计算贴近度排序。适合希望减少主观干预、评价对象较多、指标有客观数据的排序问题。常见论文表述为：
\begin{tcolorbox}
为降低主观赋权对评价结果的影响，本文采用熵权法依据各指标数据离散程度确定客观权重，并将该权重引入 TOPSIS 模型，计算各评价对象与正、负理想解的距离及相对贴近度，从而实现综合排序。
\end{tcolorbox}

\subsection{AHP-熵权组合赋权}
AHP 体现专家经验，熵权体现数据差异，可采用线性组合：
\[
w_j=\alpha a_j+(1-\alpha)b_j,
\]
其中 $a_j$ 为 AHP 主观权重，$b_j$ 为熵权，$0\le \alpha\le 1$。也可用乘法归一化：
\[
w_j=\frac{a_jb_j}{\sum_{j=1}^m a_jb_j}.
\]

\subsection{熵权-FCE}
熵权法确定指标权重，FCE 处理等级隶属关系。适合结果需要表达为“优、良、中、差”等等级，并且权重希望尽量客观的场景。

\subsection{熵权-FCE-TOPSIS 或模糊 TOPSIS}
当评价既有模糊性，又需要排序，可以先用隶属函数或模糊数刻画指标，再用 TOPSIS 计算与理想模糊解的距离；也可以先用 FCE 得到综合得分，再对多对象排序。此类模型适合环境评价、风险评价、供应商选择、项目优选等问题。

\section{建模步骤规范：可直接写进论文}
\subsection{Step1：明确评价类问题}
说明评价对象、评价目标、评价范围、评价期和数据来源。若目标是排序，优先考虑 TOPSIS；若目标是等级评价，优先考虑 FCE；若目标是客观赋权，则引入熵权法。

\subsection{Step2：建立指标体系}
构建指标集 $C=\{C_1,\ldots,C_m\}$，注明每个指标的含义、单位和属性。指标过多时可用相关分析、PCA 或专家筛选减少冗余。

\subsection{Step3：建立原始数据矩阵}
将所有评价对象和指标整理成矩阵 $X=(x_{ij})$。论文中应给出原始数据表或说明数据来源。

\subsection{Step4：正向化与数据检查}
对最小型、中间型、区间型指标进行正向化。检查缺失值、异常值、零方差指标和高度相关指标。

\subsection{Step5：标准化或归一化}
若用于 TOPSIS，可采用向量归一化或极差标准化；若用于熵权法，需要进一步构造比例矩阵；若用于 FCE，需要根据阈值、专家投票或隶属函数构造隶属度矩阵。

\subsection{Step6：确定权重}
权重可以来自 AHP、熵权法、等权法、PCA、CRITIC 或组合赋权。论文中应说明选择理由：主观经验强时用 AHP，客观数据充分时用熵权，二者都重要时用组合赋权。

\subsection{Step7：综合合成}
\begin{itemize}
  \item TOPSIS：计算正负理想解、距离和贴近度；
  \item FCE：计算 $B=WR$，根据最大隶属度或综合得分判定等级；
  \item 熵权综合评价：计算 $S_i=\sum_j w_j r_{ij}$ 或将熵权代入 TOPSIS/FCE。
\end{itemize}

\subsection{Step8：结论评价与敏感性分析}
给出排序或等级后，还要解释关键指标、短板指标和改进方向。若条件允许，可改变权重、标准化方式或删除某指标，检查结论是否稳定。

\section{选择方法的口诀}
\begin{longtable}{p{5cm}p{8.5cm}}
\toprule
研究需求 & 建议方法 \\
\midrule
只需要客观确定指标权重 & 熵权法 \\
需要多个对象排序 & TOPSIS 或熵权-TOPSIS \\
需要判断对象属于“优良中差” & FCE \\
既要等级，又要客观权重 & 熵权-FCE \\
既有专家经验又有客观数据 & AHP-熵权组合赋权 \\
既有模糊语言评价又要排序 & 模糊 TOPSIS 或 FCE 后赋分排序 \\
指标太多、相关性强 & PCA/因子分析先降维，再接 TOPSIS/FCE/熵权 \\
\bottomrule
\end{longtable}

\section{结论评价写法模板}
\subsection{统一结论模板}
\begin{tcolorbox}
本文首先构建了由 $m$ 个指标组成的综合评价指标体系，并对原始数据进行正向化和标准化处理，消除了指标方向和量纲差异。在权重确定方面，采用（AHP/熵权法/组合赋权）得到各指标权重；在综合评价方面，采用（TOPSIS/FCE/熵权综合评价）得到各评价对象的综合得分、贴近度或隶属度向量。结果表明，评价对象 $A_k$ 综合表现最好，主要原因是其在 $C_p$、$C_q$ 等关键指标上具有优势；评价对象 $A_l$ 得分较低，短板主要集中在 $C_r$ 指标。敏感性分析显示，在权重或标准化方式小幅变化时，主要排序基本保持稳定，说明模型结论具有一定稳健性。
\end{tcolorbox}

\subsection{TOPSIS 结论模板}
\begin{tcolorbox}
由 TOPSIS 计算结果可知，各评价对象相对贴近度分别为 $S_1,S_2,\ldots,S_n$。其中 $S_k$ 最大，说明对象 $A_k$ 距离正理想解最近且距离负理想解较远，综合评价最优；$S_l$ 最小，说明其综合表现与理想状态差距较大。结合指标权重可知，影响排序的主要指标为 \ldots。
\end{tcolorbox}

\subsection{FCE 结论模板}
\begin{tcolorbox}
由模糊综合评价可得综合评价向量 $B=(b_1,b_2,\ldots,b_s)$。其中 $b_k$ 最大，根据最大隶属度原则，评价对象综合等级为 $v_k$。若采用评语赋分法，综合得分为 $S=BQ^T$，说明该对象整体处于 \ldots 水平。进一步分析一级指标评价向量可知，优势因素为 \ldots，短板因素为 \ldots。
\end{tcolorbox}

\subsection{熵权法结论模板}
\begin{tcolorbox}
熵权结果显示，指标 $C_j$ 权重最高，说明该指标在当前样本中差异程度较大，对区分评价对象贡献较强；指标 $C_t$ 权重最低，说明该指标在各对象之间变化较小，区分作用有限。需要注意的是，熵权反映的是当前数据集中的区分贡献，不等同于现实意义上的绝对重要性。
\end{tcolorbox}

\section{可拓展性与较新的组合思路}
\subsection{熵权 TOPSIS}
先用熵权法确定客观权重，再用 TOPSIS 排序。它的优点是减少专家主观性，适合区域发展、环境质量、工程方案、供应商选择等多对象排序问题。

\subsection{熵权模糊综合评价}
用熵权法确定权重，用 FCE 处理等级边界。它适合水环境安全、生态安全、风险等级、教学质量等带有“优良中差”评语集的问题。

\subsection{AHP-熵权耦合}
主观权重和客观权重相结合，可避免完全专家赋权过主观，也避免纯数据赋权忽略专业经验。

\subsection{CRITIC 与熵权法对比}
熵权法重视指标离散程度；CRITIC 同时考虑指标波动性和指标之间的冲突性。若指标之间相关性很强，CRITIC 或 PCA 可以作为熵权法的对照或改进。

\subsection{模糊 TOPSIS、直觉模糊 TOPSIS}
当评价值不是精确数，而是语言评价、区间数、三角模糊数或直觉模糊数时，可以用模糊 TOPSIS。它将“模糊性”和“理想解距离”结合起来，适合专家判断不确定的决策问题。

\subsection{标准化敏感性分析}
标准化方法会影响熵权和 TOPSIS 排序，因此较严谨的论文可以比较极差标准化、向量标准化、和式归一化等方法下的权重和排序变化。

\section{网页论文实例与可以借鉴的写法}
\begin{enumerate}
  \item \textbf{TOPSIS 思路。} TOPSIS 的典型表述是：最优方案应同时距离正理想解最近、距离负理想解最远；因此适合多指标方案排序。
  \item \textbf{FCE 理论基础。} Zadeh 的模糊集理论把集合隶属关系从 0/1 扩展到 $[0,1]$，这为“优、良、中、差”等模糊等级评价提供了基础。
  \item \textbf{熵权理论基础。} Shannon 信息熵用于度量不确定性；熵权法借用这一思想，把指标在样本中的离散程度转化为客观权重。
  \item \textbf{熵权-FCE 实例。} 有研究将熵权法与模糊综合评价结合，用于饮用水源地水环境安全评价，说明“客观权重+等级隶属度”适合环境安全类问题。
  \item \textbf{熵权-TOPSIS 实例。} 有研究用熵权 TOPSIS 评价中国省域绿色发展水平，说明“客观权重+理想解排序”适合区域综合发展评价。
  \item \textbf{方法改进实例。} 近年来不少研究把 AHP、熵权、TOPSIS、FCE、模糊数、PCA、CRITIC 等组合使用，目的都是在主观经验、客观数据、模糊边界和排序稳定性之间取得平衡。
\end{enumerate}

\section{统一 Python 代码模板}
下面代码给出共同的数据处理、熵权、TOPSIS 和单层 FCE。实际建模时只需要替换原始数据矩阵、指标类型、权重或隶属度矩阵。

\begin{lstlisting}[language=Python]
import numpy as np

EPS = 1e-12

def positive_transform(X, types, targets=None, intervals=None):
    """Convert all criteria to benefit type: larger is better."""
    X = np.asarray(X, dtype=float)
    Z = X.copy()
    targets = targets or {}
    intervals = intervals or {}
    for j, t in enumerate(types):
        col = X[:, j]
        if t == "max":
            Z[:, j] = col
        elif t == "min":
            Z[:, j] = np.max(col) - col
        elif t == "mid":
            a = targets[j]
            d = np.abs(col - a)
            Z[:, j] = 1.0 if np.max(d) < EPS else 1 - d / np.max(d)
        elif t == "interval":
            a, b = intervals[j]
            d = np.zeros_like(col)
            d[col < a] = a - col[col < a]
            d[col > b] = col[col > b] - b
            Z[:, j] = 1.0 if np.max(d) < EPS else 1 - d / np.max(d)
        else:
            raise ValueError(f"Unknown criterion type: {t}")
    return Z

def minmax_standardize(Z):
    """Scale each column into [0, 1]. This does not make sums equal to 1."""
    Z = np.asarray(Z, dtype=float)
    zmin, zmax = np.min(Z, axis=0), np.max(Z, axis=0)
    denom = zmax - zmin
    R = np.zeros_like(Z)
    ok = denom > EPS
    R[:, ok] = (Z[:, ok] - zmin[ok]) / denom[ok]
    return R

def vector_normalize(Z):
    """Vector normalization, often used in TOPSIS."""
    Z = np.asarray(Z, dtype=float)
    norm = np.sqrt(np.sum(Z ** 2, axis=0))
    return Z / np.where(norm < EPS, 1.0, norm)

def entropy_weight(R):
    """Calculate entropy values and entropy weights."""
    R = np.asarray(R, dtype=float)
    m, n = R.shape
    P = (R + EPS) / np.sum(R + EPS, axis=0, keepdims=True)
    e = -np.sum(P * np.log(P), axis=0) / np.log(m)
    g = 1 - e
    w = np.ones(n) / n if np.sum(g) < EPS else g / np.sum(g)
    return e, g, w

def topsis(X, types, weights=None):
    """Entropy-weighted TOPSIS by default."""
    Z = positive_transform(X, types)
    R = vector_normalize(Z)
    if weights is None:
        _, _, weights = entropy_weight(minmax_standardize(Z))
    weights = np.asarray(weights, dtype=float)
    weights = weights / weights.sum()
    V = R * weights
    ideal_best = np.max(V, axis=0)
    ideal_worst = np.min(V, axis=0)
    d_best = np.sqrt(np.sum((V - ideal_best) ** 2, axis=1))
    d_worst = np.sqrt(np.sum((V - ideal_worst) ** 2, axis=1))
    closeness = d_worst / (d_best + d_worst + EPS)
    return closeness, weights

def fce(weights, membership_matrix, grade_scores=None):
    """Single-level fuzzy comprehensive evaluation."""
    W = np.asarray(weights, dtype=float)
    W = W / W.sum()
    R = np.asarray(membership_matrix, dtype=float)
    B = W @ R
    B_norm = B / B.sum() if B.sum() > EPS else B
    result = {"B": B_norm, "grade_index": int(np.argmax(B_norm))}
    if grade_scores is not None:
        result["score"] = float(B_norm @ np.asarray(grade_scores, dtype=float))
    return result

if __name__ == "__main__":
    X = np.array([
        [85, 12, 70, 0.80],
        [90, 16, 65, 0.75],
        [78, 10, 80, 0.88],
    ])
    types = ["max", "min", "max", "max"]
    c, w = topsis(X, types)
    print("Entropy weights:", np.round(w, 4))
    print("TOPSIS closeness:", np.round(c, 4))

    R = np.array([
        [0.70, 0.20, 0.10, 0.00, 0.00],
        [0.30, 0.40, 0.20, 0.10, 0.00],
        [0.50, 0.30, 0.20, 0.00, 0.00],
        [0.20, 0.40, 0.30, 0.10, 0.00],
    ])
    res = fce([0.25, 0.25, 0.25, 0.25], R, [100, 80, 60, 40, 20])
    print("FCE vector:", np.round(res["B"], 4), "score:", round(res["score"], 2))
\end{lstlisting}

\section{常见错误清单}
\begin{enumerate}
  \item 把标准化和归一化混为一谈。标准化后落在 $[0,1]$，不代表和为 1。
  \item 忘记正向化，导致“越小越好”的指标被当成“越大越好”。
  \item 把 TOPSIS 的贴近度当成概率。$S_i$ 是相对接近程度，不是事件发生概率。
  \item 把 FCE 的隶属度当成概率。隶属度表示属于某等级的程度，只有频率法构造时才具有近似统计频率解释。
  \item 把熵权解释为现实中的绝对重要性。熵权反映的是当前数据差异贡献。
  \item 只给最终得分，不解释权重、关键指标和短板指标。
  \item 没有说明最大值、最小值来自理论标准还是当前样本。
  \item 指标高度重复却直接赋权，导致某类信息被重复计算。
\end{enumerate}

\section{最后总总结}
TOPSIS、FCE 和熵权法可以统一理解为多指标评价问题的不同模块：
\begin{itemize}
  \item \textbf{熵权法}解决“权重怎么客观确定”；
  \item \textbf{TOPSIS}解决“多个对象如何按照接近理想状态排序”；
  \item \textbf{FCE}解决“评价等级边界模糊时如何分级”；
  \item \textbf{AHP}解决“专家经验如何进入权重”；
  \item \textbf{PCA、CRITIC、Bootstrap、敏感性分析}解决“指标冗余、权重偏差和结论稳健性”。
\end{itemize}

所以，数学建模论文中可以把这类方法的本质写成：\textbf{建立指标体系，整理数据矩阵，完成同向化与无量纲化，确定权重，再通过合适的综合评价算子将多指标信息转化为排序、等级和改进建议。}

\section*{参考资料与网页论文实例}
\begin{enumerate}
  \item Shannon, C. E. A Mathematical Theory of Communication. Bell System Technical Journal, 1948. \url{https://people.math.harvard.edu/~ctm/home/text/others/shannon/entropy/entropy.pdf}
  \item Zadeh, L. A. Fuzzy Sets. Information and Control, 1965. \url{https://www.sciencedirect.com/science/article/pii/S001999586590241X}
  \item ScienceDirect Topics. Technique for Order Preference by Similarity to Ideal Solution. \url{https://www.sciencedirect.com/topics/engineering/technique-for-order-preference-by-similarity-to-ideal-solution}
  \item Ding, X. et al. Fuzzy Comprehensive Assessment Method Based on the Entropy Weight Method and Its Application in Water Environmental Safety Evaluation. Water, 2017. \url{https://www.mdpi.com/2073-4441/9/5/329}
  \item Zhao, H. et al. A Fuzzy Comprehensive Evaluation Method Based on AHP and Entropy for a Landslide Susceptibility Map. Entropy, 2017. \url{https://www.mdpi.com/1099-4300/19/8/396}
  \item Ma, X. et al. Green Development Evaluation Based on Entropy Weight TOPSIS. International Journal of Environmental Research and Public Health, 2023. \url{https://www.mdpi.com/1660-4601/20/3/1707}
  \item Roszkowska, E. et al. Impact of Normalization on Entropy-Based Weights. Entropy, 2024. \url{https://www.mdpi.com/1099-4300/26/5/365}
\end{enumerate}

\end{document}
